\section{Introduction}\label{sec:intro}

Strong earthquakes in steep, tectonically active terrain trigger thousands to tens of thousands of landslides within seconds. These landslides cause a large share of earthquake fatalities, cut transport corridors and start sediment cascades that persist for years \citep{keefer1984,fan2019}. The 2016 M\textsubscript{w}~7.8 Kaik\={o}ura earthquake ruptured more than twenty faults \citep{hamling2017} and triggered roughly 30\,000 landslides over about \SI{10000}{\square\kilo\metre} of northeastern South Island \citep{massey2018,massey2020update}. A future Alpine Fault earthquake is expected to produce a comparable or larger landslide response in the Southern Alps \citep{robinson2016}, and response planning for that scenario depends on landslide forecasts for an event that has not happened \citep{orchiston2018}. The same problem recurs around the Pacific Rim, from the Andes and Central America to Japan, Alaska, Papua New Guinea and Sulawesi.

In the first hours after an earthquake, no inventory exists for the new event. Landslide models used in emergency response therefore have to be \emph{transferred}: they are fitted to past earthquakes and applied to a new one. Operational near-real-time systems do this with logistic regression on ShakeMap ground motion and global terrain and geology layers \citep{nowicki2014,nowickijessee2018,allstadt2022}. Physics-based alternatives use Newmark rigid sliding-block displacement with generic strength parameters \citep{newmark1965,jibson2000,jibson2007}. In parallel, most landslide studies now use machine-learning (ML) classifiers such as gradient-boosted trees, because these usually give the highest discrimination when trained and tested within one region \citep{merghadi2020,reichenbach2018}. Discrimination inside a single event says little about transfer. Spatially autocorrelated samples inflate random cross-validation scores \citep{brenning2012,roberts2017}, and high scores can coexist with geomorphically implausible maps \citep{steger2016}. A flexible learner can also fit event-specific artefacts, such as inventory mapping protocol, local ShakeMap error or lithology, that do not carry over to the next earthquake.

Cross-event transfer has been tested before. Global models are evaluated on events withheld from fitting \citep{nowicki2014,nowickijessee2018,tanyas2019slope}, and fuzzy and regression models have been carried from one earthquake to another \citep{kritikos2015,li2013lushan}. \citet{lombardo2020chrono} framed this as chrono-validation, and \citet{allstadt2018} evaluated candidate global models against the Kaik\={o}ura inventory. These studies leave four questions open, and they are the gaps this paper addresses.

\begin{enumerate}
\item No controlled, like-for-like comparison exists between physics-based, global-style logistic and ML models under leave-one-event-out validation with identical inputs. It is therefore unknown whether the advantage of ML learners survives transfer to a new earthquake, or whether adding physical knowledge as monotonic constraints \citep{karpatne2017,reichstein2019} improves it.
\item Transfer studies rarely separate \emph{discrimination}, the ranking of hillslopes, from \emph{calibration}, the agreement of predicted probabilities with observed landslide frequencies. Calibration is the property most likely to fail in a new population \citep{vancalster2019}, and probability-based products depend on it \citep{allstadt2022}.
\item The value of the first few landslides mapped from post-event imagery \citep{williams2018,handwerger2022} for updating a transferred model has not been quantified, even though clinical prediction research provides established updating methods \citep{steyerberg2004,janssen2008}.
\item Explainable-AI attributions \citep{lundberg2017,lundberg2020} are increasingly read as evidence of process understanding in landslide studies \citep{dahal2023xai}. Whether these attributions stay the same from one earthquake to the next, which they must if a model has learned process, has not been tested.
\end{enumerate}

We address these gaps with nine coseismic landslide inventories from around the Pacific Rim: Kaik\={o}ura 2016 (New Zealand), Northridge 1994 and Denali 2002 (USA), Niigata-Chuetsu 2004 and Iwate-Miyagi 2008 (Japan), Ays\'en 2007 (Chile), Cinchona 2009 (Costa Rica), Porgera 2018 (Papua New Guinea) and Palu 2018 (Indonesia). We use only globally and rapidly available inputs: Copernicus GLO-30 terrain, which is derived from TanDEM-X radar interferometry \citep{copdem,guth2021}, and USGS ShakeMap ground motion \citep{worden2018}. With these inputs we (i) benchmark five models of increasing flexibility in a leave-one-event-out (LOEO) design, including a physics-constrained gradient-boosting model (PC-GBM) whose response to shaking, slope and Newmark displacement is forced to be monotone; (ii) separate transfer error into discrimination and calibration components; (iii) quantify how many rapidly mapped image tiles are needed to recalibrate a transferred model; and (iv) test whether TreeSHAP attributions are stable across events. Kaik\={o}ura is used throughout as the target case for Aotearoa New Zealand.
